\chapter{Estudio previo}\label{cap:planificación}

\section{Introducción}
Breve introducción al capítulo

\section{ArangoDB}
Es una base de datos NoSQL nativa multimodelo, lo que significa que permite trabajar con documentos, grafos y sistemas clave-valor utilizando un único motor central y un mismo lenguaje de consultas llamado AQL (ArangoDB Query Language).

\section{Arquitectura Multimodelo}
En el mundo de las bases de datos, normalmente estás atado a un solo formato: tablas (relacional), documentos tipo JSON (como MongoDB), o grafos (como Neo4j).

La principal ventaja de la arquitectura multimodelo nativa de ArangoDB es que combina varios modelos en un mismo motor central, permitiendo acceder a ellos con el mismo lenguaje de consultas (AQL). Los tres modelos que soporta son:
\begin{itemize}
\item \textbf{Documentos:} Para guardar información jerárquica y flexible.
\item \textbf{Grafos:} Para definir relaciones complejas y direcciones entre los datos.
\item \textbf{Clave-Valor:} Para accesos de lectura ultrarrápidos mediante un identificador único.
\end{itemize}

\section{Modelado de Grafos y AQL}
En una base de datos tradicional, buscar conexiones requiere cruzar tablas mediante operaciones costosas. En un modelo de grafos, las relaciones son importantes y trazarlas es directo e inmediato.
AQL (ArangoDB Query Language) es el lenguaje que permite recorrer esta red de datos. .

\section{Integración}
